% ============================================================
% Notes: Integral Equations (A&W Ch. 16.3--16.4)
% Topics: types, kernels, Neumann series, separable kernels,
%         eigenvalues/eigenfunctions of homogeneous Fredholm II
% ============================================================

\section{Integral equations}

\subsection{General form}
An integral equation is an equation where the unknown function $\phi(x)$ appears under
an integral sign. A standard linear form is
\begin{equation}
\phi(x) = f(x) + \lambda \int_a^b K(x,t)\,\phi(t)\,dt,
\end{equation}
where
\begin{itemize}
\item $f(x)$ is known (driving term),
\item $K(x,t)$ is the kernel,
\item $\lambda$ is a constant parameter,
\item $\phi(x)$ is the unknown.
\end{itemize}

% ------------------------------------------------------------

\section{Types of integral equations}

\subsection{Fredholm vs Volterra}
\paragraph{Fredholm.}
Limits are constants:
\begin{equation}
\phi(x) = f(x) + \lambda \int_a^b K(x,t)\phi(t)\,dt.
\end{equation}

\paragraph{Volterra.}
One limit depends on $x$ (typical causal structure):
\begin{equation}
\phi(x) = f(x) + \lambda \int_a^x K(x,t)\phi(t)\,dt.
\end{equation}

\subsection{First kind vs second kind}
\paragraph{First kind.}
Unknown appears only inside the integral:
\begin{equation}
f(x) = \int_a^b K(x,t)\phi(t)\,dt.
\end{equation}
These are often ill-conditioned / non-unique.

\paragraph{Second kind.}
Unknown appears both outside and inside:
\begin{equation}
\phi(x) = f(x) + \lambda \int_a^b K(x,t)\phi(t)\,dt.
\end{equation}
This is typically better posed and supports iterative solution methods.

% ------------------------------------------------------------

\section{Neumann series (successive approximations)}

\subsection{Setup (Fredholm second kind)}
Consider the Fredholm equation of the second kind
\begin{equation}
\phi(x) = f(x) + \lambda \int_a^b K(x,t)\phi(t)\,dt.
\end{equation}
Neumann's idea: iterate the equation by substituting approximations into the integral.
A natural first guess is
\begin{equation}
\phi_0(x) = f(x).
\end{equation}
Then define recursively
\begin{equation}
\phi_{n+1}(x) = f(x) + \lambda \int_a^b K(x,t)\phi_n(t)\,dt.
\end{equation}

\subsection{Series form}
Iterating generates a power series in $\lambda$:
\begin{equation}
\phi(x) = \sum_{i=0}^{\infty} \lambda^i u_i(x),
\end{equation}
with
\begin{align}
u_0(x) &= f(x),\\
u_1(x) &= \int_a^b K(x,t_1)f(t_1)\,dt_1,\\
u_2(x) &= \int_a^b\!\!\int_a^b K(x,t_1)K(t_1,t_2)f(t_2)\,dt_2\,dt_1,\\
u_n(x) &= \int_a^b\cdots\int_a^b
K(x,t_1)K(t_1,t_2)\cdots K(t_{n-1},t_n)\,f(t_n)\,dt_n\cdots dt_1.
\end{align}
This is exactly the Neumann series structure developed in Arfken \& Weber. :contentReference[oaicite:2]{index=2}

\subsection{Kernel iterates (iterated kernels)}
Define the iterated kernels
\begin{align}
K_1(x,t) &= K(x,t),\\
K_{n+1}(x,t) &= \int_a^b K(x,z)\,K_n(z,t)\,dz.
\end{align}
Then the Neumann series may be written compactly as
\begin{equation}
\phi(x)=f(x)+\lambda\int_a^b K_1(x,t)f(t)\,dt
+\lambda^2\int_a^b K_2(x,t)f(t)\,dt+\cdots.
\end{equation}

\subsection{Convergence idea}
The Neumann series converges if the integral operator is ``small enough'' (in operator norm),
or more precisely (as emphasized in the text) convergence is controlled by the eigenvalues
of the corresponding homogeneous equation. :contentReference[oaicite:3]{index=3}

\subsection{Volterra case (important simplification)}
For Volterra equations of the second kind,
\begin{equation}
\phi(x)=f(x)+\lambda\int_a^x K(x,t)\phi(t)\,dt,
\end{equation}
the Neumann series converges for all $\lambda$ provided $K$ is square integrable. :contentReference[oaicite:4]{index=4}

% ------------------------------------------------------------

\section{Separable (degenerate) kernels}

\subsection{Definition}
A kernel is \textbf{separable} (degenerate) if it can be written as a finite sum
\begin{equation}
K(x,t)=\sum_{j=1}^n M_j(x)\,N_j(t),
\end{equation}
as in A\&W. :contentReference[oaicite:5]{index=5}
Example:
\begin{equation}
\cos(t-x)=\cos t\cos x+\sin t\sin x.
\end{equation}

\subsection{Fredholm second kind with separable kernel}
Substitute into
\begin{equation}
\phi(x)=f(x)+\lambda\int_a^b K(x,t)\phi(t)\,dt
\end{equation}
to obtain
\begin{equation}
\phi(x)=f(x)+\lambda\sum_{j=1}^n M_j(x)\int_a^b N_j(t)\phi(t)\,dt.
\end{equation}
Define constants
\begin{equation}
c_j:=\int_a^b N_j(t)\phi(t)\,dt.
\end{equation}
Then the solution has the form
\begin{equation}
\phi(x)=f(x)+\lambda\sum_{j=1}^n c_j M_j(x).
\end{equation}

\subsection{Reduction to linear algebra}
Multiply by $N_i(x)$ and integrate to obtain a finite linear system:
\begin{equation}
c_i=b_i+\lambda\sum_{j=1}^n a_{ij}c_j,
\end{equation}
where
\begin{equation}
b_i=\int_a^b N_i(x)f(x)\,dx,
\qquad
a_{ij}=\int_a^b N_i(x)M_j(x)\,dx.
\end{equation}
Matrix form:
\begin{equation}
b=(1-\lambda A)c,
\qquad
c=(1-\lambda A)^{-1}b.
\end{equation}
This makes explicit the correspondence between integral equations and matrix equations. :contentReference[oaicite:6]{index=6}

\subsection{Homogeneous case and eigenvalues}
For the homogeneous equation $f(x)=0$, we have $b=0$ and
\begin{equation}
(1-\lambda A)c=0.
\end{equation}
Nontrivial solutions exist only if
\begin{equation}
\boxed{\det(1-\lambda A)=0.}
\end{equation}
This is the \textbf{secular equation} giving eigenvalues $\lambda$. :contentReference[oaicite:7]{index=7}

% ------------------------------------------------------------

\section{Homogeneous Fredholm equation of the second kind}

\subsection{Eigenvalue formulation}
The homogeneous Fredholm equation of the second kind is
\begin{equation}
\boxed{
\phi(x)=\lambda\int_a^b K(x,t)\phi(t)\,dt.
}
\end{equation}
A nontrivial solution $\phi(x)\neq 0$ exists only for special values $\lambda=\lambda_n$.
These are the \textbf{eigenvalues}, and the corresponding solutions $\phi_n(x)$ are
the \textbf{eigenfunctions}. :contentReference[oaicite:8]{index=8}

\subsection{Symmetric kernels and Hilbert--Schmidt theory}
Assume
\begin{equation}
K(x,t)=K(t,x),\qquad K\in\mathbb{R}.
\end{equation}
Then (Hilbert--Schmidt theory):
\begin{itemize}
\item eigenvalues $\lambda_n$ are real,
\item eigenfunctions corresponding to distinct eigenvalues are orthogonal,
\item eigenfunctions form a complete set in a suitable sense.
\end{itemize}
This parallels Hermitian matrices and Sturm--Liouville ODEs. :contentReference[oaicite:9]{index=9}

\subsection{Orthogonality of eigenfunctions}
Let $\phi_i,\phi_j$ satisfy
\[
\phi_i(x)=\lambda_i\int_a^b K(x,t)\phi_i(t)\,dt,
\qquad
\phi_j(x)=\lambda_j\int_a^b K(x,t)\phi_j(t)\,dt.
\]
For a symmetric kernel, one obtains
\begin{equation}
(\lambda_j-\lambda_i)\int_a^b \phi_i(x)\phi_j(x)\,dx = 0.
\end{equation}
Thus if $\lambda_i\neq \lambda_j$,
\begin{equation}
\boxed{
\int_a^b \phi_i(x)\phi_j(x)\,dx = 0.
}
\end{equation}
This is the Fredholm integral-equation analog of Sturm--Liouville orthogonality. :contentReference[oaicite:10]{index=10}

\subsection{Normalization and degeneracy}
One may normalize eigenfunctions so that
\begin{equation}
\int_a^b \phi_i(x)\phi_j(x)\,dx=\delta_{ij}.
\end{equation}
If an eigenvalue is degenerate, the corresponding eigenfunctions may be orthogonalized
by Gram--Schmidt. :contentReference[oaicite:11]{index=11}

\subsection{Reality of eigenvalues}
For Hermitian kernels $K(x,t)=K^*(t,x)$ (general complex case), eigenvalues are real. :contentReference[oaicite:12]{index=12}

% ------------------------------------------------------------

\section{Kernel expansion in eigenfunctions (bilinear expansion)}

For symmetric kernels, the kernel itself can be expanded in eigenfunctions:
\begin{equation}
\boxed{
K(x,t)=\sum_{n=1}^{\infty}\frac{\phi_n(x)\phi_n(t)}{\lambda_n},
\qquad (0\ \text{not an eigenvalue}).
}
\end{equation}
This is the bilinear eigenfunction expansion stated in A\&W. :contentReference[oaicite:13]{index=13}

% ------------------------------------------------------------

\section{High-yield summary}

\begin{itemize}
\item Fredholm: fixed limits; Volterra: variable limit.
\item First kind: $f=\int K\phi$ (often non-unique/ill-conditioned).
\item Second kind: $\phi=f+\lambda\int K\phi$ (iterative methods work).
\item Neumann series: iterate substitution to get a power series in $\lambda$:
\[
\phi=\sum_{n=0}^\infty \lambda^n u_n.
\]
\item Separable (degenerate) kernels reduce the integral equation to a finite matrix equation:
\[
(1-\lambda A)c=b.
\]
\item Homogeneous Fredholm II: $\phi=\lambda\int K\phi$ has nontrivial solutions only for
eigenvalues $\lambda_n$.
\item For symmetric kernels: eigenvalues real, eigenfunctions orthogonal, completeness holds.
\item Kernel eigenfunction expansion:
\[
K(x,t)=\sum_n \frac{\phi_n(x)\phi_n(t)}{\lambda_n}.
\]
\end{itemize}
